% !TEX root = ../../main-jos.tex
\section{引言}

微服务架构把一个大系统拆成多个独立服务，带来了更好的可扩展性和开发灵活性。容器编排、Sidecar 代理和服务治理等技术的成熟，使云原生应用的大规模部署成为现实\cite{burns2016patterns,burns2016borg,jamshidi2018microservices,richardson2018microservices,feng2020microservice,wuhuayao2020microdev}。拆分也带来一个棘手问题：日志高度分散。每个服务实例各自输出日志，格式五花八门，总量随并发请求线性增长。传统做法是``全量采集''，先把日志统统收上来再说。这在中小规模场景下还能接受，一旦集群达到数十甚至数百节点，问题就暴露了：节点 CPU 和内存被持续占用，南北向带宽成为瓶颈，存储与索引成本快速攀升，而且大量低价值日志反而把真正有用的故障信息给淹没了\cite{liao2016logs,limingshu2019logmgmt,he2020logsurvey,usman2022observability}。

目前的研究主要集中在日志存储压缩、采样过滤、故障诊断和 eBPF 无侵入采集等方向，但很少有人从\textbf{网关流量感知}的角度出发，动态驱动各节点的应用日志定向采集。包航宇等\cite{bao2023aiops}总结了智能运维的实践现状与标准化框架，贾统等\cite{jiatong2020logdiag}系统梳理了基于日志的分布式故障诊断技术。本文的思路不同，把关注点前移到\textbf{采集前端}——网关是南北向流量的入口，天然掌握着 URL、状态码、响应时延等结构化信号。利用这些信号，不需要侵入业务代码，就能识别出哪些请求是高价值的（比如出错的、响应慢的），然后指导下游节点只采集与这些请求相关的应用日志，做到``只采必要日志''。

\textbf{本文关注的核心问题是}：在不侵入业务代码、不改变微服务调用链的前提下，能否利用网关层的流量异常信号，从采集前端大幅降低日志入库量，同时保留异常和慢请求的高价值上下文？围绕这个核心问题，本文拆解出四个科学子问题：（1）如何从采集前端实现大幅减量，同时保留故障诊断所需的关键上下文？（2）定向采集会不会给 Sidecar 带来不可接受的资源开销？（3）如何在不侵入业务代码的前提下，利用网关流量信号动态识别高价值请求，并确保采集策略在网关预筛选、节点定向采集和中心二次过滤三层之间保持语义一致，形成端到端的定向闭环？（4）如何在 Sidecar 资源受限环境下实现可靠的日志缓存与传输，避免因中心暂时不可达或网络拥塞导致的数据丢失？

\textbf{本文提出以下科学假设}：（H1）网关流量驱动的关注清单能使定向采集模式下的 Loki 入库量相比同架构全量采集显著下降，同时在规则级保留异常与慢请求的关键上下文。（H2）定向采集的 Sidecar 资源开销（CPU、内存）不高于同架构全量采集。（H3）网关访问日志中的 URL、状态码和响应时延足以动态生成覆盖异常与慢请求模式的关注清单；通过三次策略转换，可以维持跨层策略语义一致，形成端到端的采集闭环。（H4）固定缓存块与指数退避机制能够在低资源占用下提供有界重试和传输缓冲能力。后文实验部分以 RQ1--RQ4 对上述假设逐一验证。

基于上述问题与假设，本文提出\textbf{分布式定向日志采集框架}（Distributed Directed Log Collection Framework）。框架分为三层：\textbf{网关预筛选、节点定向采集、中心二次过滤}，三层之间通过三次策略转换保持语义一致。和 Promtail/Fluent Bit ``全量推送后在存储端压缩''的传统方案相比，本文直接在采集前端就把 Loki 入库量控制在百条量级（八节点实测详见 §6.3），从源头降低了网络和存储开销。主要架构贡献如下：

（1）\textbf{三层协同定向采集架构}。网关节点和微服务节点均以 Sidecar 方式部署，中心负责策略生成与二次过滤，节点负责本地匹配、缓存与可靠上传，构建了``流量感知—定向采集—语义一致入库''的完整闭环。这个架构将日志减量逻辑从传统的后端治理前置到采集前端，补上了 AIOps 数据链路中``采集前端减量''这一缺失环节。

（2）\textbf{动态关注度评分模型（DASM）}。提出以频次、错误率、延迟严重度和热度趋势四个归一化因子加权求和的动态评分机制（式~\ref{eq:dasm}），结合指数时间衰减（式~\ref{eq:decay}）和负载自适应权重调节，实时生成高价值 URL 模式清单并下发至各节点（算法~\ref{alg:attention}，排序实现复杂度上界为 $O(N\log N)$，最小堆实现为 $O(N\log K)$，$K \ll N$）。和简单的频次$\times$严重度加权相比，DASM 能够感知模式的时序变化趋势并自动适应不同负载特征，同时满足评分有界性（定理~\ref{thm:bound}）和退化兼容性（命题~\ref{prop:degrade}）。

（3）\textbf{定向策略三次转换机制}。通过统一的 URL 泛化函数和清单版本控制，确保同一关注清单在网关预筛选、节点定向采集和中心二次过滤三层之间的语义一致性，避免跨层策略漂移（图~\ref{fig:transform}）。

（4）\textbf{资源受限环境下的可靠传输架构}。在 Sidecar 中引入固定缓存块（环形队列+双约束）和压力感知指数退避机制（式~\ref{eq:backoff}），辅以 BoltDB 本地兜底，在资源压力下提供有界内存占用和有界重试能力。

（5）\textbf{真实部署验证}。在 8 微服务 DSB-Lite 真实部署环境下，$n{=}5$ 配对实验显示 DASM 减量 $25.9\% \pm 0.9\%$（$p<0.001$），且不引入额外延迟开销。

（6）\textbf{与 OpenTelemetry 尾采样的工程对比}。在相同负载下对比 OTel Collector v0.103 + Tail Sampling，验证了 DASM 在\emph{高价值端点保留率}和\emph{决策延迟}两个维度上的优势。

基于 Go、gRPC、Redis 和 Grafana Loki 构建了可复现原型和多进程模拟器。在八节点集群上完成 180\,s、$n{=}10$ 的重复对比实验（95\% CI 自助法估计）、算法微基准测试与系统级组件消融，并补充了 DSB-Lite 真实部署验证、OpenTelemetry 尾采样对比、1--1000 节点扩展性外推。结果显示，定向采集将 Loki 入库量从 4388$\pm$1 条降至 72 条（降幅 98.4\%，策略过滤独立贡献 67.8\%）。DSB-Lite 真实部署上减量 $25.9\% \pm 0.9\%$（$p<0.001$）。十六/三十二节点复核验证了定向模式在更大规模下依然能保持低入库量和低 Agent CPU 开销。

本文第 2 节综述相关工作；第 3 节介绍系统总体设计；第 4 节阐述关键算法；第 5 节描述实现细节；第 6 节给出实验与分析；第 7 节总结全文。
